\documentclass[11pt]{article}
\usepackage{amsmath,amssymb,amsthm,hyperref}
\newtheorem{theorem}{Theorem}
\newtheorem{proposition}[theorem]{Proposition}
\title{A Rank Obstruction for Split-Product Cancellation}
\author{}
\date{November 23, 2026}
\begin{document}
\maketitle

\paragraph{Scope.}
This note audits a proposed extension of the reciprocal cycle identities.
It does not prove the requested universal head bound. The polynomial
identities below are identities over the real polynomial ring, not merely
identities on the Boolean cube. The rank theorem used below is an external
theorem; this note does not claim a Lean formalization of it.

\section{The exact construction being tested}

Let $u_1,\ldots,u_m,v_1,\ldots,v_m$ be real affine forms in $x$, and set
\[
 F(s,x)=\prod_{j=1}^m(s+u_j(x)),\qquad
 G(s,x)=\prod_{j=1}^m(s+v_j(x)).
\]
Suppose $F-G=R(x)$ as a polynomial identity, with $R$ independent of $s$.
Differentiation gives $F_s=G_s$ and hence
\[
 \sum_{j=1}^m\frac1{s+u_j(x)}-
 \sum_{j=1}^m\frac1{s+v_j(x)}
 =\frac{F_s}{F}-\frac{G_s}{G}
 =-\frac{R(x)F_s(s,x)}{F(s,x)G(s,x)}.
\]
At every point where all the displayed affine factors are positive,
$F_s=\sum_j\prod_{k\ne j}(s+u_k)>0$. Therefore the score has the
strict sign of $-R$ wherever $R\ne0$. Substitution of an affine
$s=C(x)$ preserves both affine denominators and this identity. The
substitution must separately ensure the manuscript's denominator
admissibility. The displayed score uses $2m$ heads before cancellations.

Thus a tempting construction is to arrange a prescribed product of affine
sign tests as $F-G$, and then use its logarithmic derivative difference.
The next result bounds the number of independent directions in this
construction, independently of $m$.

\section{Rank from a three-term polynomial identity}

We use Saxena and Seshadhri, \emph{From Sylvester--Gallai Configurations
to Rank Bounds: Improved Black-Box Identity Test for Depth-3 Circuits},
ECCC TR10-013, revision 1,
\url{https://eccc.weizmann.ac.il/report/2010/013/revision/1/download/}.
Their Theorem 5 states that a simple, minimal, identically zero real
depth-three circuit with $k$ product terms has linear-factor span of
dimension less than $3k^2$. Here simple means that the product terms
have no common linear factor, and minimal means that no nonempty proper
subsum vanishes.

\begin{proposition}[Split-product differences have bounded factor rank]
Let $F,G,R$ be nonzero homogeneous real polynomials of the same positive
degree, each a scalar times a product of real linear forms, with
$F-G=R$. Divide all three polynomials by their greatest common divisor.
The linear factors of the resulting three polynomials span a vector
space of dimension at most $26$.
\end{proposition}
\begin{proof}
Write the quotients as $F_0,G_0,R_0$. They are nonzero, have the same
degree, and satisfy $F_0-G_0-R_0=0$. The three product terms have no
common linear factor. No one-term subsum is zero. If a two-term subsum
were zero, the remaining nonzero term would have to be zero by the
identity. Thus the circuit is simple and minimal. The cited theorem
with $k=3$ gives rank less than $27$, hence rank at most $26$.
If the quotients have degree zero, there are no remaining linear factors
and the conclusion is immediate without applying the theorem.
\end{proof}

The rank bound must not be silently replaced by three. For independent
variables $a,b,c,d$ there is the identity
\[
 (a^2-b^2)(c^2-d^2)-(a^2-c^2)(b^2-d^2)
 +(a^2-d^2)(b^2-c^2)=0.
\]
Every displayed quadratic factor splits over the reals. The products
are pairwise coprime, and their linear factors span all four variables:
for example $a+b,a-b,c+d,c-d$ are independent. The identity follows
by expanding with $A=a^2,B=b^2,C=c^2,D=d^2$. It is consistent with
the bound $26$ used here.

For affine products, homogenize every factor to degree one in a new
variable $z$. If the difference has smaller degree, homogenize it to
the common product degree, including the required power of $z$.
If the affine difference is a product of real affine forms, this
homogenized difference is still completely split over the reals.
The proposition therefore applies to the homogenized identity.
Dehomogenization and affine substitutions are linear maps on the
coefficient vectors and cannot increase their span dimension.

Common factors of $F,G$ also divide $F-G$. They cancel identically in
$\partial_s F/F-\partial_s G/G$, because the logarithmic derivative
of their common product appears in both summands. Consequently the
remaining denominators in this score have slope vectors spanning at
most $26$ dimensions. This assertion includes multiplicities.

\paragraph{A hypothesis that is not being assumed.}
The preceding argument does not require an irreducible generic fiber.
Pereira--Yuzvinsky, Theorem 2.1(3), gives a sharper bound through a
linear pullback from projective four-space when that additional
hypothesis holds. Coprimality alone must not be used to infer it:
$F=X^2$ and $G=Y^2$ are coprime, whereas their generic pencil member
splits over the algebraic closure. The real circuit-rank theorem
avoids that inference entirely.

\section{The cost of combining independent cancellation blocks}

The following counting statement permits arbitrary affine numerators,
so it enlarges the class of logarithmic derivative constructions.

\begin{theorem}[Bounded-rank blocks]
Fix constants $r\ge1$ and $C>0$. Let $N=2^n$ and
$1\le K\le H\le CN/n^2$. Consider scores with at most $H$ positive
affine denominators. Suppose the heads can be partitioned into at most
$K$ groups so that, within each group, the denominator slope vectors
span a space of dimension at most $r$. The groups, their spanning
vectors, and all coefficients may depend on the Boolean target.
If $K/H\to0$, the number of strict Boolean sign tables realized by
this class is $2^{o(N)}$.
\end{theorem}
\begin{proof}
Normalize each denominator by its positive value at $x=0$. Absorbing
each numerator's constant multiple of the denominator into the bias
leaves $k_1=1+nH$ numerator parameters. For a fixed group assignment
$g:[H]\to[K]$, write
\[
 B_h(x)=1+\sum_{a=1}^r t_{ha}\sum_{i=1}^n v_{g(h),a,i}x_i.
\]
The denominator block has $q=rnK+rH$ parameters. Every allowed
group of slope vectors admits such a factorization, including groups
of rank less than $r$. The cleared numerator at each cube vertex has
degree at most one in the numerator block and at most $2H$ in the
denominator block. With $k=k_1+q$, the two-block sign-count lemma
in \texttt{mechanism-audit.tex} gives, for sufficiently large $n$,
\[
 K^H\left(\frac{32768N}{k}\right)^k2^{k_1}(2H)^q
\]
as an upper bound on the number of strict patterns. The factor $K^H$
counts all assignments, and $k\le N$ holds because
$k=O_{C,r}(N/n)$. Applying the sign count on all real parameters only
enlarges the family obtained by requiring positive admissible
denominators.

The base-two logarithm is bounded by
\[
 O_{C,r}(N\log n/n)+O_r(n^2K+nH).
\]
Indeed $H\log_2 K=O_C(N/n)$,
$k\log_2(32768N/k)=O_{C,r}(N\log n/n)$, and
$q\log_2(2H)=O_r(n^2K+nH)$.
Finally $n^2K\le CN(K/H)=o(N)$ and $nH=O_C(N/n)$.
This proves the claim. Zero-numerator constant-denominator heads
allow padding smaller representations without increasing group rank.
\end{proof}

\paragraph{Consequence for the proposed construction.}
Apply the theorem with $r=26$ to any sum of split-product logarithmic
derivative blocks of the kind just proved. At the requested head
budget, a family using $o(H)$ such blocks cannot be universal.
Increasing the degree of a small number of split-product identities
therefore cannot supply the missing factor of $n$.
The result does not exclude a number of small blocks comparable to
$H$, differences that do not split over the reals, or identities valid
only modulo the Boolean ideal. It is not a lower bound for the full
manuscript model.

For the more specific monic construction $F-G=R(x)$ with $R$
independent of $s$, \texttt{ideal-moment-rank-audit.tex} subsequently
proves the sharper rank bound three without requiring $R$ to split.
That argument uses moment annihilation and an elementary
Cayley--Bacharach interpolation bound. The rank-$26$ statement above
still applies to the broader class of three fully split products
without the common monic auxiliary-variable hypothesis.

\paragraph{Counting audit.}
The proof deliberately uses two parameter blocks. Replacing them by
one block per head and multiplying only the individual degree factors
would omit the dependence of multiblock sign bounds on the number of
blocks. In particular, that shortcut would incorrectly strengthen the
bound even when every group contains just one arbitrary denominator.
No such shortcut is used above.

\end{document}
